% Section 3: Methodology
\section{Methodology}
\label{sec:methodology}

Building on our observation that the capability-LC alignment question is inherently a partial
correlation problem, we design a multi-layer statistical analysis establishing existence (P1),
mechanism (P2), and population-level confirmation (P3).

\subsection{Dataset}

We use the publicly available AlpacaEval 2.0 leaderboard CSV (accessed August 2026),
containing pairwise preference results for 223 diverse LLMs compared against a GPT-4o reference
on 805 instruction-following prompts. The three variables of interest are:

\begin{itemize}
  \item \textbf{\texttt{win\_rate}}: Fraction of pairwise comparisons where human annotators
    prefer the model (range: ${\sim}3$--$80\%$)
  \item \textbf{\texttt{LC\_winrate}}: Win rate after GLM-based length debiasing
    \citep{Dubois2024LengthControlled} (range: ${\sim}2$--$95\%$)
  \item \textbf{\texttt{avg\_length}}: Mean response token count per model
    (range: ${\sim}400$--$2200$ tokens)
\end{itemize}

All 223 rows have complete data; no imputation was required.

\textbf{Verbosity Separability Check}: VIF = 1.764 for both \texttt{win\_rate} and
\texttt{avg\_length} (threshold: VIF $<$ 5.0), confirming OLS coefficients and partial
correlations are interpretable (Figure~\ref{fig:vif}).

\subsection{Primary Analysis: Partial Correlation (H-E1)}

We seek $\rho(\texttt{win\_rate}, \texttt{LC\_winrate} \mid \texttt{avg\_length})$ --- the
Spearman partial correlation after controlling for \texttt{avg\_length}.

\textbf{Implementation}: \texttt{pingouin.partial\_corr(method='spearman', alternative='greater')}
with covariate \texttt{avg\_length}. Pre-registered threshold: $|r_{\mathrm{partial}}| \geq 0.15$,
$p < 0.05$ (one-tailed; directional hypothesis $\rho > 0$). Bootstrap 95\% CI: 1000 resamples,
\texttt{random\_state=42}.

\subsection{Mechanism Analysis: Standardized OLS (H-M1)}

\textbf{Design}: \texttt{LC\_winrate} $\sim$ \texttt{win\_rate\_std} + \texttt{avg\_length\_std}
with \texttt{sklearn.StandardScaler} preprocessing and \texttt{statsmodels.OLS}. Dominance
criterion: $|\beta_{\mathrm{win}}| > |\beta_{\mathrm{len}}|$. Verbosity-only baseline
(\texttt{LC\_winrate} $\sim$ \texttt{avg\_length\_std}) for $R^2$ comparison.

\subsection{FWL-Inspired Robustness Verification (H-M2)}

\textbf{Design}: Regress \texttt{avg\_length} from both \texttt{win\_rate} and
\texttt{LC\_winrate} (OLS); compute Spearman $\rho$ on residual pairs. Delta
$= |\rho_{\mathrm{Pingouin}} - \rho_{\mathrm{resid}}| < 0.02$ confirms estimator consistency.
\textit{Note}: The Frisch--Waugh--Lovell theorem guarantees exact algebraic equality between
partial OLS coefficients and residualized OLS estimates (Pearson). For Spearman partial
correlation, residualized ranks provide an independent estimator that approximates this
principle; convergence within 1.1pp is empirical evidence of robustness, not a theorem
guarantee.

\subsection{Population-Level Confirmation: Kruskal-Wallis + Dunn (H-M3, H-C1)}

\textbf{Design}: \texttt{win\_rate} quartiles (Q1--Q4; $N = 56, 56, 55, 56$).
\texttt{scipy.stats.kruskal} on both \texttt{LC\_winrate} (H-C1) and $\Delta$ (H-M3).
Post-hoc: \texttt{scikit\_posthocs.posthoc\_dunn} with Bonferroni correction. Effect size:
$\varepsilon^2 = (H - k + 1) / (N - k)$. Bootstrap CI on Dunn Q1 vs Q4 $p$ (1000 resamples).

\subsection{Sub-Hypothesis Structure}

\begin{table}[h]
\centering
\caption{Sub-hypothesis summary}
\label{tab:hypotheses}
\begin{tabular}{llll}
\toprule
Hypothesis & Type & Gate & Test \\
\midrule
H-E1 & Existence & MUST\_WORK & Partial corr $> 0.15$ \\
H-M1 & Mechanism & MUST\_WORK & $|\beta_{\mathrm{win}}| > |\beta_{\mathrm{len}}|$ in standardized OLS \\
H-M2 & Mechanism & SHOULD\_WORK & FWL delta $< 0.02$ \\
H-M3 & Mechanism & SHOULD\_WORK & KW on $\Delta$ across quartiles, $p < 0.05$ \\
H-C1 & Condition & SHOULD\_WORK & KW on LC\_winrate AND Dunn Q1 vs Q4, both $p < 0.05$ \\
\bottomrule
\end{tabular}
\end{table}
